\section{Introduction}
\label{sec:intro}

Advertising systems predict several types of conversion at once, and those types differ sharply
in how long they take to observe. A website view or an add-to-cart is recorded within minutes of the
click; a purchase, an app install or a submitted lead may take a day or a week. Each objective
is defined with an \emph{attribution window}, the elapsed time within which a conversion is
credited to the click. Conventionally that window also determines the \emph{join delay}, the age at which the click's
training record is written to the stream: the pipeline waits until the label is knowable, then
emits the record. This conversion has two disadvantages: A model for a slow objective trains only on data
at least an attribution window old, so it never sees fresh data during online training. And a system that ranks several
objectives jointly must either maintain one pipeline per objective or delay all of them until
the slowest window closes, imposing the staleness of the slowest funnel on the fastest.

Neither those two disadvantages are intrinsic to conversion rate modeling. While the attribution window is fixed by the business definition of a
conversion, the join delay is a pipeline choice. We ask what happens when the two are set
independently, and specifically whether an objective with attribution window of days can be
trained from records joined within minutes.

The obstacle is that an early-joined record is wrong: every conversion still in
flight when the record is written is recorded as a negative. A substantial literature corrects
this censoring
bias~\cite{chapelle2014modeling,yasui2020feedback,yang2021capturing,chen2022asymptotically,ding2026delayed},
but most corrections revise the record afterwards, by duplicating the click when its outcome
lands, by relabelling it, or by adding prediction heads at future horizons. A streaming
pipeline that appends each record once and trains on it once can do none of these. We state
that restriction as a \emph{single-visit} constraint.

Our idea is this censoring bias can be mitigated by the prediction from a model that has seen trained on the matured outcomes in the conventional way, i.e. \emph{a teacher model}.
Such a teacher model is often readily available in production, where large ranking models are already trained
across objectives on matured data~\cite{meta2025gem}. The contributions of this work are: (1) a separation of the join delay from the attribution
window, with the causality and single-visit constraint a deployable pipeline must satisfy. After
auditing \nScreened{} published methods, \nFail{} methods violate single-visit by
writing more than 1 record per click; (2) evidence that writing the record early and
replacing only its observed negatives with teacher model predictions beats waiting for the matured label; (3) the finding that at small
served capacities the hard label can be dropped altogether, since a training pipeline consisting only of
the teacher's estimations as labels can beat the conventional training pipeline with delayed join same as attribution windows, which can benefit early stage, such as retrieval, where model complexity is low.
